\clearpage
\VerdierPhysicalPageAnchor{0230}%
\begingroup\small\setlength{\parskip}{0pt}\setlength{\emergencystretch}{1.5em}
\setlength{\abovedisplayskip}{4pt plus 2pt minus 2pt}%
\setlength{\belowdisplayskip}{4pt plus 2pt minus 2pt}%

\VerdierSemanticAnchor{ch4d-4-1-6-cont-2}
\noindent puis :

\VerdierSemanticAnchor{ch4d-4-1-6-6}
\begin{equation}
\int_n f\bigl((p,q),(p',q')\bigr)
=h\bigl((p,q),(p',q')\bigr).
\tag{4.1.6.6}
\end{equation}

Soit de plus

\VerdierSemanticAnchor{ch4d-4-1-6-7}
\begin{equation}
{}^{m}\!d(p,q,r)\colon{}^{m}\!Y(q,r)
\longrightarrow{}^{m}\!Y(p,q)[1]
\tag{4.1.6.7}
\end{equation}

le morphisme composé :

\VerdierSemanticAnchor{ch4d-4-1-6-8}
\begin{equation}
\begin{aligned}
{}^{m}\!Y(q,r)
&\xrightarrow{\int_n g(p,q,r)}
\int_n T_m\bigl({}^{m}\!X(p,q)\bigr)
\xrightarrow{(*)}
\left(\int_n{}^{m}\!X(p,q)\right)[1]\\
&=Y(p,q)[1].
\end{aligned}
\tag{4.1.6.8}
\end{equation}

où le morphisme $(*)$ est l'isomorphisme de commutation aux translations du
foncteur complexe simple associé (chap.~I, 2.3.2.1).

On obtient ainsi une filtration du complexe simple $Y=\int_n X$ :

\VerdierSemanticAnchor{ch4d-4-1-6-9}
\begin{equation}
0={}^{m}\!Y(-\infty)\subset\cdots\subset{}^{m}\!Y(q)
\subset\cdots\subset{}^{m}\!Y(+\infty)=Y
\tag{4.1.6.9}
\end{equation}

qu'on appelle la \emph{filtration définie par le $m$-ième degré de $X$}.

On peut, par le procédé rappelé en (4.1.1), associer au complexe filtré
(4.1.6.9) un objet spectral à valeurs dans $K(\mathcal A)$ noté :

\VerdierSemanticAnchor{ch4d-4-1-6-10}
\begin{equation}
\bigl(\mathcal Y(p,q),\delta(p,q,r)\bigr).
\tag{4.1.6.10}
\end{equation}

Enfin, abusivement, nous désignerons dans la proposition ci-après par la
même notation un morphisme de $\operatorname{comp}(\mathcal A)$ et son
image dans $K(\mathcal A)$.

\VerdierSemanticAnchor{ch4d-4-1-7}
\medskip\noindent\textbf{Proposition 4.1.7.}\quad

\VerdierSemanticAnchor{ch4d-4-1-7-a}
\noindent\textit{a) Les objets ${}^{m}\!Y(p,q)$ (4.1.6.5), les morphismes
$h\bigl((p,q),(p',q')\bigr)$ (4.1.6.6) et les morphismes
${}^{m}\!d(p,q,r)$ (4.1.6.7) définissent un objet spectral à valeurs dans
$K(\mathcal A)$ qui dépend fonctoriellement du complexe $n$-uple $X$.}

\noindent\textit{Cet objet spectral est appelé l'objet spectral à valeurs
dans $K(\mathcal A)$ défini par le complexe $X$ filtré par son $m$-ième
degré.}

\VerdierSemanticAnchor{ch4d-4-1-7-b}
\noindent\textit{b) Il existe un isomorphisme canonique (i.e.\ fonctoriel par
rapport à l'objet $X$ de $\operatorname{comp}^n(\mathcal A)$) entre les
objets spectraux $\bigl(\mathcal Y(p,q),\delta(p,q,r)\bigr)$ (4.1.6.10) et
$\bigl({}^{m}\!Y(p,q),{}^{m}\!d(p,q,r)\bigr)$.}
\endgroup
